% Shared macros for the Lee "Introduction to Riemannian Manifolds" (2nd ed.) blueprint.
% Kept KaTeX-compatible (no fancy packages) to match the leanblueprint web build.

\newcommand{\R}{\mathbb{R}}
\newcommand{\C}{\mathbb{C}}
\newcommand{\Z}{\mathbb{Z}}
\newcommand{\N}{\mathbb{N}}
\newcommand{\Rn}{\R^n}
\newcommand{\Sph}[1]{\mathbb{S}^{#1}} % round sphere, e.g. \Sph{n}
\newcommand{\Hyp}[1]{\mathbb{H}^{#1}} % hyperbolic space, e.g. \Hyp{n}

% Tangent/cotangent spaces and bundles
\newcommand{\T}{T} % tangent bundle / tangent space, e.g. \T_p M, \T M
\newcommand{\Tp}[1]{T_{#1}} % T_p, T_q, ...
\newcommand{\Ts}{T^*} % cotangent bundle
\newcommand{\Tsp}[1]{T^*_{#1}}

% Metric, connection, curvature
\newcommand{\metric}{g}
\newcommand{\conn}{\nabla} % Levi-Civita / general connection
\newcommand{\Christoffel}[3]{\Gamma_{#1#2}^{#3}}
\newcommand{\Riem}{R} % (3,1) curvature tensor
\newcommand{\Rm}{\mathrm{Rm}} % (0,4) curvature tensor
\newcommand{\Ric}{\mathrm{Ric}}
\newcommand{\Scal}{S} % scalar curvature
\newcommand{\Sec}{\mathrm{Sec}} % sectional curvature
\newcommand{\Weyl}{W} % Weyl tensor
\newcommand{\grad}{\operatorname{grad}}
\newcommand{\Hess}{\operatorname{Hess}}
\newcommand{\dive}{\operatorname{div}}
\newcommand{\Exp}{\operatorname{Exp}}
\newcommand{\Ker}{\operatorname{Ker}}
\newcommand{\sgn}{\operatorname{sgn}}
\newcommand{\supp}{\operatorname{supp}}
\newcommand{\cof}{\operatorname{cof}}
% Kulkarni-Nomizu product (circled wedge), e.g. h \owedge k
\newcommand{\owedge}{\mathbin{\text{\textcircled{\raisebox{-0.9pt}{$\scriptstyle\wedge$}}}}}

\newcommand{\inner}[2]{\langle #1, #2 \rangle}
\newcommand{\norm}[1]{\lvert #1 \rvert}
\newcommand{\dist}{d} % Riemannian distance function

% Book cross-reference macro: \dcref{chN:N.M} points at an item's own numbering
% in Lee's book (e.g. "Theorem 2.5" in Chapter 2 becomes \dcref{ch2:2.5}), kept
% as a plain no-op placeholder here; the print/web leanblueprint build is free to
% turn it into a hyperlink or margin note later.
\newcommand{\dcref}[1]{}

% Sign flip of a bilinear form b along a vector u: b - 2 b(u,.) (x) b(u,.)
\newcommand{\flip}{\operatorname{fl}}
